
\chapter{Sensitivity Analysis for the Average Causal Effect with Unmeasured Confounding}
\chaptermark{Sensitivity Analysis for Average Causal Effect}\label{chapter::sensitivity-ACE}


 
 
 
 
在给定观测协变量的条件下，Cornfield-type sensitivity analysis 最适合二元结局的 risk ratio 尺度。虽然 \citet{ding2016sensitivity} 也为 average causal effect 提出了 Cornfield-type sensitivity analysis 方法，但这些方法还不够一般，实际应用也不够方便。下面给出一种更直接的方法：它以 potential outcomes 的条件期望为基础来做 sensitivity analysis。这个方法的优势在于，它可以把 average causal effect 的常用估计量纳入 sensitivity analysis 框架中。
这个想法已经出现在 \citet{robins1999association} 和 \citet{scharfstein1999adjusting} 的早期工作中。本章基于 \citet{luding2023flexible} 最近的表述。
 
 
 
 
 这个方法与为 average potential outcomes 推导 worst-case bounds 的思想密切相关。我先回顾较简单的 bounds 思想，然后把它扩展到 sensitivity analysis。



\section{Introduction}


回顾 observational study 的标准设定：$\{  Z_i, X_i,  Y_i(1), Y_i(0) \}_{i=1}^n \iidsim  \{ Z, X, Y(1), Y(0) \}$，我们关注 average causal effect
$$
\tau = E\{ Y(1) - Y(0) \}.
$$
它可以分解为
\begin{eqnarray*}
\tau &=& \left[  E(Y\mid Z=1) \pr(Z=1) + E\{ Y(1)\mid Z=0 \} \pr(Z=0) \right]  \\
&&- \left[   E\{ Y(0)\mid Z=1\} \pr(Z=1) + E(Y\mid Z=0)  \pr(Z=0)  \right] . 
\end{eqnarray*}
因此，根本困难在于估计 counterfactual means
$$
E\{ Y(1)\mid Z=0 \}  ,\quad \quad E\{ Y(0)\mid Z=1\}.
$$
一般来说，估计它们有两种极端策略。



我们已经在\cref{part::observational-studies}讨论过第一种策略，它依赖 ignorability。
假设
\begin{eqnarray*}
E\{ Y(1)\mid Z=1, X \} &=& E\{ Y(1)\mid Z=0, X \} ,\\ 
E\{ Y(0)\mid Z=1, X \} &=& E\{ Y(0)\mid Z=0, X \}, 
\end{eqnarray*}
我们可以用可观测量识别 counterfactual means：
$$
E\{ Y(1)\mid Z=0 \}  
%= E\left[ E\{ Y(1)\mid Z=0, X \}  \mid Z=0 \right] 
%= E\left[ E\{ Y(1)\mid Z=1, X \}  \mid Z=0\right] 
= E\left\{  E( Y\mid Z=1, X )  \mid Z=0\right\}
$$
同理，
$$
E\{ Y(0)\mid Z=1\} = E\left\{  E( Y\mid Z=0, X )  \mid Z=1\right\} . 
$$



下一节中的第二种策略除了假设结局被限制在 $\underline{y}$ 和 $\overline{y}$ 之间以外，不作任何假设。
对于二元结局，这个假设很自然，此时 $\underline{y} = 0$ 且 $\overline{y} = 1$。
在这个假设下，两个 counterfactual means 也都被限制在 $\underline{y}$ 和 $\overline{y}$ 之间，从而给出 $\tau$ 的 worst-case bounds。
\cref{table::manski-bounds} 展示了基本思想，下面的\cref{sec::manski-bounds}会更详细地回顾这一策略。


 


 



\begin{table}
\centering
\caption{带有有界结局 $ [\underline{y}, \overline{y}]$ 的 Science Table，其中 $\underline{y}$ 和 $\overline{y}$ 是两个常数}\label{table::manski-bounds}
\begin{tabular}{ccccccc}
\hline 
$Z$ & $Y(1)$ & $Y(0)$ &  Lower $Y(1)$ & Upper $Y(1)$ & Lower $Y(0)$ & Upper $Y(0)$ \\
\hline 
$1$ & $Y_1(1)$ & ? & $Y_1(1)$  & $Y_1(1)$ & $\underline{y}$ & $\overline{y}$ \\
$\vdots$&$\vdots$&$\vdots$&$\vdots$&$\vdots$&$\vdots$&$\vdots$   \\
$1$ & $Y_{n_1}(1)$ & ? &  $Y_{n_1}(1)$  & $Y_{n_1}(1)$ & $\underline{y}$ & $\overline{y}$ \\
\hline 
$0$ & ? & $Y_{n_1+1}(0)$ & $\underline{y}$ & $\overline{y}$& $Y_{n_1+1}(0)$ & $Y_{n_1+1}(0)$ \\ 
$\vdots$&$\vdots$&$\vdots$&$\vdots$&$\vdots$&$\vdots$&$\vdots$ \\ 
$0$ & ? & $Y_{n}(0)$ & $\underline{y}$ & $\overline{y}$& $Y_{n}(0)$ & $Y_{n}(0)$ \\
\hline 
\end{tabular}
\end{table}


\section{Manski-type worse-case bounds on the average causal effect without assumptions}\label{subsection::manski-bounds}
\label{sec::manski-bounds}



假设结局被限制在 $\underline{y}$ 和 $\overline{y}$ 之间。
由分解式
$$
E\{  Y(1) \} = E\{  Y(1) \mid Z=1\} \pr(Z=1)  +  E\{  Y(1) \mid Z=0\} \pr(Z=0)  ,
$$
可知 $E\{  Y(1) \}$ 的下界为
$$
E\{  Y \mid Z=1\} \pr(Z=1)  +   \underline{y} \pr(Z=0) 
$$
上界为
$$
E\{  Y \mid Z=1\} \pr(Z=1)  +  \overline{y} \pr(Z=0).
$$
类似地，由分解式
$$
E\{  Y(0) \} = E\{  Y(0) \mid Z=1\} \pr(Z=1)  +  E\{  Y(0) \mid Z=0\} \pr(Z=0)  ,
$$
可知 $E\{  Y(0) \}$ 的下界为
$$
\underline{y}\pr(Z=1)  +  E\{  Y \mid Z=0\} \pr(Z=0) 
$$
上界为
$$
 \overline{y} \pr(Z=1)  +  E\{  Y \mid Z=0\} \pr(Z=0). 
$$
合并这些 bounds，可以得到 average causal effect $\tau = E\{  Y(1) \} - E\{  Y(0) \}$ 的下界
$$
E\{  Y \mid Z=1\} \pr(Z=1)  +   \underline{y} \pr(Z=0) - \overline{y} \pr(Z=1)  -  E\{  Y \mid Z=0\} \pr(Z=0) 
$$
以及上界
$$
 E\{  Y \mid Z=1\} \pr(Z=1)  +  \overline{y} \pr(Z=0) - \underline{y}\pr(Z=1)  -   E\{  Y \mid Z=0\} \pr(Z=0). 
$$ 
这个区间的长度为 $\overline{y}-\underline{y}$。这些 bounds 信息量不强，但仍好于先验 bounds $[\underline{y}-\overline{y}, \overline{y}-\underline{y}]$，后者长度为 $2(\overline{y}-\underline{y})$。在没有进一步假设时，观测数据分布不能唯一决定 $\tau$。在这种情形下，我们说 $\tau$ 是 {\it partially identified}，其形式定义如下。


\begin{definition}
[partial identification]\label{def::partial-identification}
如果观测数据分布可以与 $\theta$ 的多个取值相容，则参数 $\theta$ 是 partially identified。
\end{definition}


比较\cref{def::nonparametric-identification}和\cref{def::partial-identification}。如果参数 $\theta$ 由观测数据分布唯一决定，那么它是 identifiable；否则，它只是 partially identifiable。因此，在 ignorability assumption 下，$\tau$ 是 identifiable；没有 ignorability assumption 时，它只是 partially identifiable。




\citet{cochran1953} 在有 missing data 的调查中使用过 worst-case bounds 的思想，但后来放弃了这一想法，因为它经常给出非常保守的结果。类似地，上面关于 $\tau$ 的 worst-case bounds 从实践角度看往往不够有趣，因为它们常常包含 $0$。此外，这一策略不适用于结局无界的设定。



Manski 将这一思想用于 causal inference \citep{manski1990nonparametric} 以及许多其他 econometric models \citep{manski2003partial}。
用最少假设来约束 causal parameters 的思想，在与其他 qualitative assumptions 结合时非常有力。\citet{manski2003partial} 综述了许多策略。
例如，我们可能相信 treatment 不会伤害任何 unit，因此 monotonicity assumption 成立：
$
Y(1)\geq Y(0) . 
$
那么 $\tau$ 的下界为零，而上界保持不变。
另一类假设是
$Z = I\{  Y(1) \geq Y(0) \}$，也就是说，treatment selection 基于 latent potential outcomes 之间的差异。这个假设也可以改进 $\tau$ 的 bounds。
关于这一方法的更详细讨论超出了本书范围。

 
 


 
\section{Sensitivity analysis for the average causal effect}


 
第一种策略是乐观的：它假设在给定观测协变量后，potential outcomes 在 treatment group 和 control group 之间没有差异。第二种策略是悲观的：它完全不根据观测数据来推断 counterfactual means。下面的策略介于二者之间。



\subsection{Identification formulas}

定义
\begin{eqnarray*}
\frac{ E\{ Y(1)\mid Z=1, X \} }{ E\{ Y(1)\mid Z=0, X \} } = \varepsilon_1(X), \\
\frac{ E\{ Y(0)\mid Z=1, X \} }{ E\{ Y(0)\mid Z=0, X \} } = \varepsilon_0(X) ,
\end{eqnarray*}
它们就是 sensitivity parameters。为简单起见，我们还可以进一步假设它们是不依赖于 $X$ 的常数。在实际中，我们需要固定这些参数，或者让它们在预先指定的范围内变化。回顾 $\mu_1(X) = E(Y\mid Z=1, X)$ 和 $\mu_0(X) = E(Y\mid Z=0, X)$ 分别是 treatment 和 control 下观测结局的条件均值函数。我们可以按如下方式识别两个 counterfactual means 和 average causal effect。

\begin{theorem}
\label{thm::outcome-reg-sensitivity}
若 $\varepsilon_1(X)$ 和 $ \varepsilon_0(X) $ 已知，则有
\begin{eqnarray*}
E\{ Y(1)\mid Z=0 \}   &=& E\left\{\mu_1(X)   / \varepsilon_1(X)  \mid Z=0\right\},\\ 
E\{ Y(0)\mid Z=1 \}   &=& E\left\{\mu_0(X)    \varepsilon_0(X)  \mid Z=1\right\}
\end{eqnarray*} 
因此
\begin{eqnarray}
\tau & = & E\{ ZY + (1-Z) \mu_1(X)   / \varepsilon_1(X)\}  \nonumber  \\
 && - E\{ Z \mu_0(X)    \varepsilon_0(X) + (1-Z)Y\}  \label{eq::predictive}\\
& = & E\{ Z \mu_1(X) + (1-Z) \mu_1(X)   / \varepsilon_1(X)\}  \nonumber \\
&&- E\{ Z \mu_0(X)    \varepsilon_0(X) + (1-Z) \mu_0(X)\}. \label{eq::projective}
\end{eqnarray}
\end{theorem}


 


\cref{thm::outcome-reg-sensitivity} 的证明留给\cref{hw::proof-theorem-1-sa}。
在拟合 outcome model 后，\cref{eq::predictive} 和 \cref{eq::projective} 启发了下面两个用于估计 $\tau$ 的 predictive estimator 和 projective estimator：
\begin{eqnarray*}
\hat\tau^{\textup{pred}} &=&\left\{  n^{-1}\sumn Z_i Y_i + n^{-1} \sumn (1-Z_i) \hat{\mu}_1(X_i)   / \varepsilon_1(X_i)  \right\} \\
&&- \left\{ n^{-1}  \sumn Z_i \hat{\mu}_0(X_i)    \varepsilon_0(X_i) + n^{-1} \sumn (1-Z_i) Y_i \right\},
\end{eqnarray*}
以及
\begin{eqnarray*}
\hat{\tau}^{\textup{proj}}&=&\left\{ n^{-1}\sumn Z_{i}\hat{\mu}_{1}(X_{i})+n^{-1}\sumn(1-Z_{i})\hat{\mu}_{1}(X_{i})/\varepsilon_{1}(X_{i})\right\} \\
&&-\left\{ n^{-1}\sumn Z_{i}\hat{\mu}_{0}(X_{i})\varepsilon_{0}(X_{i})+n^{-1}\sumn(1-Z_{i})\hat{\mu}_{0}(X_{i})\right\}.
\end{eqnarray*} 
``predictive'' 和 ``projective'' 这两个术语来自 survey sampling 文献 \citep{firth1998robust, ding2018causal}；也见\cref{sec::understand-lin-predict}。估计量 $\hat\tau^{\textup{pred}} $ 和 $\hat{\tau}^{\textup{proj}}$ 略有不同：前者在观测结局可用时使用观测结局，而后者用拟合值替代观测结局。





更有意思的是，我们还可以用 inverse probability weighting 公式来识别 $\tau$。

\begin{theorem}\label{thm::ipw-sensitivity}
若 $\varepsilon_1(X)$ 和 $ \varepsilon_0(X) $ 已知，则有
\begin{eqnarray*} 
E\{ Y(1) \} = 
E\left\{ w_1(X) \frac{Z}{e(X)}  Y  \right\}  ,\quad
 E\{Y(0)\} = E\left\{ w_0(X) \frac{1-Z}{1-e(X)} Y \right\},
\end{eqnarray*}
其中
$$
w_1(X) = e(X) + \{1-e(X)\}/ \varepsilon_1(X) ,\quad
w_0(X) = e(X)  \varepsilon_0(X) + 1- e(X). 
$$
\end{theorem}


 
 

\cref{thm::ipw-sensitivity} 的证明留给\cref{hw::proof-theorem-2-sa}。
\cref{thm::ipw-sensitivity} 用两个额外因子 $w_1(X) $ 和 $w_0(X) $ 修正经典 IPW 公式；这两个因子同时依赖 propensity score 和 sensitivity parameters。在拟合 propensity scores 后，\cref{thm::ipw-sensitivity} 启发了下面用于估计 $\tau$ 的估计量：
\begin{eqnarray*} 
\hat\tau^{\textup{ht}} &=&  n^{-1} \sumn \frac{ \{ \hat{e}(X_i)\varepsilon_1(X_i) + 1-\hat{e}(X_i) \} Z_iY_i     }{  \varepsilon_1(X_i)  \hat e(X_i)  } \\
&&- n^{-1} \sumn  \frac{ \{  \hat e(X_i)  \varepsilon_0(X_i) + 1- \hat e(X_i)  \} (1-Z_i) Y_i }{ 1- \hat e(X_i)} 
\end{eqnarray*}
以及
\begin{eqnarray*} 
\hat\tau^{\textup{haj}} &=&   \sumn \frac{ \{ \hat{e}(X_i)\varepsilon_1(X_i) + 1-\hat{e}(X_i) \} Z_iY_i     }{  \varepsilon_1(X_i)  \hat  e(X_i)  }  \Big / \sumn \frac{Z_i}{ \hat{e}(X_i)  }   \\
&&- n^{-1} \sumn  \frac{ \{  \hat  e(X_i)  \varepsilon_0(X_i) + 1- \hat e(X_i)  \} (1-Z_i) Y_i }{ 1- \hat e(X_i)} \Big / \sumn \frac{1-Z_i}{ 1- \hat{e}(X_i)  } .
\end{eqnarray*}





更进一步，如果同时拟合 propensity score model 和 outcome models，则下面的 $\tau$ 估计量是 doubly robust 的：
\begin{eqnarray*}
\hat{\tau}^{\textup{ht}} = \hat\tau^{\textup{ht}} 
- n^{-1} \sumn \{Z_i- \hat e(X_i)\}  \left\{   \frac{  \hat  \mu_1(X_i) }{  \hat e(X_i) \varepsilon_1(X_i)  }    + \frac{  \hat  \mu_0(X_i)\varepsilon_0(X_i)  }{ 1-  \hat  e(X_i) } \right\}.
\end{eqnarray*}
也就是说，当 $\varepsilon_1(X_i) $ 和 $\varepsilon_0(X_i) $ 已知时，只要 propensity score model 或 outcome model 至少有一个设定正确，估计量 $\hat{\tau}^{\textup{dr}} $ 就是 $\tau$ 的 consistent estimator。我们可以用 bootstrap 近似上述估计量的方差。技术细节见 \citet{luding2023flexible}。




当 $\varepsilon_1(X_i)  =  \varepsilon_0(X_i) = 1$ 时，上述估计量退化为\cref{part::observational-studies}中介绍的 predictive estimator、IPW estimator 和 doubly robust estimators。



\subsection{Example}
\label{sec::eg-sensitivity-tau}


\subsubsection{R functions for sensitivity analysis}


下面的 \ri{R} 函数可以计算 sensitivity analysis 的点估计。

\begin{lstlisting}
OS_est_sa = function(z, y, x, out.family = gaussian, 
                     truncps = c(0, 1), e1 = 1, e0 = 1)
{
     ## fitted propensity score
     pscore   = glm(z ~ x, family = binomial)$fitted.values
     pscore   = pmax(truncps[1], pmin(truncps[2], pscore))
     
     ## fitted potential outcomes
     outcome1 = glm(y ~ x, weights = z, 
                    family = out.family)$fitted.values
     outcome0 = glm(y ~ x, weights = (1 - z), 
                    family = out.family)$fitted.values
     
     ## outcome regression estimator
     ace.reg  = mean(z*y) + mean((1-z)*outcome1/e1) - 
                   mean(z*outcome0*e0) - mean((1-z)*y) 
     ## IPW estimators
     w1 = pscore + (1-pscore)/e1
     w0 = pscore*e0 + (1-pscore)
     ace.ipw0 = mean(z*y*w1/pscore) - 
                   mean((1 - z)*y*w0/(1 - pscore))
     ace.ipw  = mean(z*y*w1/pscore)/mean(z/pscore) - 
                   mean((1 - z)*y*w0/(1 - pscore))/mean((1 - z)/(1 - pscore))
     ## doubly robust estimator
     aug = outcome1/pscore/e1 + outcome0*e0/(1-pscore)
     ace.dr   = ace.ipw0 + mean((z-pscore)*aug)

     return(c(ace.reg, ace.ipw0, ace.ipw, ace.dr))     
}
\end{lstlisting}

standard errors 的计算留给\cref{hw::sensitivity-se}。

\subsubsection{Revisit Example \cref{eg::chan-bmi-ATE}}


取
$$
\varepsilon_1(X) = \varepsilon_0(X) \in \{1/2, 1/1.7, 1/1.5, 1/1.3, 1, 1.3, 1.5, 1.7, 2\},
$$ 
根据下面的 \ri{R} 代码，我们得到一组 $\tau$ 的 doubly robust estimates：
\begin{lstlisting}
> nhanes_bmi = read.csv("nhanes_bmi.csv")[, -1]
> z = nhanes_bmi$School_meal
> y = nhanes_bmi$BMI
> x = as.matrix(nhanes_bmi[, -c(1, 2)])
> x = scale(x)
> 
> 
> E1 = c(1/2, 1/1.7, 1/1.5, 1/1.3, 1, 1.3, 1.5, 1.7, 2)
> E0 = c(1/2, 1/1.7, 1/1.5, 1/1.3, 1, 1.3, 1.5, 1.7, 2)
> EST = outer(E1, E0)
> ll1 = length(E1)
> ll0 = length(E0)
> for(i in 1:ll1)
+   for(j in 1:ll0)
+     EST[i, j] = OS_est_sa(z, y, x, e1 = E1[i], e0 = E0[j])[4]
\end{lstlisting}




\begin{table}[ht]
\centering
\caption{average causal effect 的 sensitivity analysis}\label{table::sensitivity-ACE}
\begin{tabular}{rrrrrrrrrr}
  \hline
 & 1/2 & 1/1.7 & 1/1.5 & 1/1.3 & 1 & 1.3 & 1.5 & 1.7 & 2 \\ 
  \hline
1/2 & 11.62 & 10.44 & 9.40 & 8.03 & 4.96 & 0.97 & -1.69 & -4.35 & -8.34 \\ 
  1/1.7 & 9.22 & 8.05 & 7.00 & 5.64 & 2.57 & -1.42 & -4.08 & -6.75 & -10.74 \\ 
  1/1.5 & 7.63 & 6.45 & 5.41 & 4.05 & 0.97 & -3.02 & -5.68 & -8.34 & -12.33 \\ 
  1/1.3 & 6.03 & 4.86 & 3.81 & 2.45 & -0.62 & -4.61 & -7.27 & -9.94 & -13.93 \\ 
  1 & 3.64 & 2.47 & 1.42 & 0.06 & -3.01 & -7.01 & -9.67 & -12.33 & -16.32 \\ 
  1.3 & 1.80 & 0.63 & -0.42 & -1.78 & -4.85 & -8.85 & -11.51 & -14.17 & -18.16 \\ 
  1.5 & 0.98 & -0.19 & -1.24 & -2.60 & -5.67 & -9.66 & -12.33 & -14.99 & -18.98 \\ 
  1.7 & 0.36 & -0.82 & -1.86 & -3.23 & -6.30 & -10.29 & -12.95 & -15.61 & -19.60 \\ 
  2 & -0.35 & -1.52 & -2.57 & -3.93 & -7.00 & -10.99 & -13.65 & -16.32 & -20.31 \\ 
   \hline
\end{tabular}
\end{table}


\cref{table::sensitivity-ACE} 给出了点估计。
估计值的符号对大于 $1$ 的 sensitivity parameters 并不敏感，但对小于 $1$ 的 sensitivity parameters 相当敏感。当 meal plan 的参与者倾向于有更高的 BMI 时（也就是说，$\varepsilon_1(X) >1$ 且 $ \varepsilon_0(X)  > 1$），meal plan 对 BMI 的 average causal effect 为负。不过，如果 meal plan 的参与者倾向于有更低的 BMI，这一结论就可能相当敏感。



\section{Homework Problems}


\homeworksolutionref{18}
\paragraph{Proof of \cref{thm::outcome-reg-sensitivity}}\label{hw::proof-theorem-1-sa}

证明\cref{thm::outcome-reg-sensitivity}。

\paragraph{Proof of \cref{thm::ipw-sensitivity}}\label{hw::proof-theorem-2-sa}

证明\cref{thm::ipw-sensitivity}。


\paragraph{Standard errors in sensitivity analysis}\label{hw::sensitivity-se}

\cref{sec::eg-sensitivity-tau} 只给出了点估计。报告对应的 bootstrap standard errors。

  


\paragraph{Sensitivity analysis for the average causal effect on the treated units $\tau_\textsc{T}$}\label{hw::att-sentivity-analysis}

本题将\cref{chapter::ATT-and-other}扩展到允许 unmeasured confounding 的情形，以估计
$$
\tau_\textsc{T} = E\{  Y(1) - Y(0)  \mid Z=1  \} 
= E(Y\mid Z=1) - E\{  Y(0)  \mid Z=1  \} .
$$
我们可以很容易地用样本矩估计 $E(Y\mid Z=1)$，即 $\hat{\mu}_{\textsc{t}1} =   \sumn Z_i Y_i / \sumn Z_i$。唯一的 counterfactual term 是 $E\{  Y(0)  \mid Z=1  \}$。因此，我们只需要 sensitivity parameter $ \varepsilon_0(X) $。当 $\varepsilon_0(X)$ 已知时，有下面两个 identification formulas。

\begin{theorem}
\label{thm::att-reg-ipw}
若 $\varepsilon_0(X) $ 已知，则有
 \begin{eqnarray*}
E\{  Y(0)  \mid Z=1  \}  &=& E\left\{  Z \mu_0(X)  \varepsilon_0(X)    \right\} /e \\
&=& E\left\{  e(X)\varepsilon_0(X) \frac{ 1-Z}{1-e(X)}  Y      \right\} /e,
\end{eqnarray*}
其中 $e = \pr(Z=1)$。
\end{theorem}


证明\cref{thm::att-reg-ipw}。

Remark:  \cref{thm::att-reg-ipw} 启发我们使用 $\hat{\tau}_{\textsc{t}}^* = \hat{\mu}_{\textsc{t}1}  - \hat{\mu}_{\textsc{t}0}^*$ 来估计 $\tau_{\textsc{t}}$，其中
 \begin{eqnarray*}
\hat{\mu}_{\textsc{t}0}^\text{reg} &=& n_1^{-1} \sumn Z_i  \varepsilon_0(X_i) \hat  \mu_0(X_i) ,   \\
\hat{\mu}_{\textsc{t}0}^\text{ht} &=& n_1^{-1} \sumn   \varepsilon_0(X_i) \hat  o(X_i) (1-Z_i) Y_i ,   \\
\hat{\mu}_{\textsc{t}0}^\text{haj} &=&  \sumn   \varepsilon_0(X_i) \hat  o(X_i) (1-Z_i) Y_i  /  \sumn   \hat  o(X_i) (1-Z_i) , 
\end{eqnarray*}
其中 $  \hat  o(X_i) = \hat  e(X_i) / \{1- \hat  e(X_i)\}$ 是给定观测协变量时 treatment 的 estimated conditional odds。
此外，我们可以为 $\tau_{\textsc{t}}$ 构造 doubly robust estimator
$\hat{\tau}_{\textsc{t}}^\text{dr} = \hat{\mu}_{\textsc{t}1}  - \hat{\mu}_{\textsc{t}0}^\text{dr}$，其中
$$
\hat{\mu}_{\textsc{t}0}^\text{dr} = 
\hat{\mu}_{\textsc{t}0}^\text{ht}  - n_1^{-1} \sumn \varepsilon_0(X_i) \frac{\hat e(X_i) - Z}{1-\hat e(X_i)}  \hat \mu_0(X_i)  . 
$$
\citet{luding2023flexible} 给出了更多细节。



\paragraph{R code}

实现\cref{hw::att-sentivity-analysis}中的估计量。
分析\cref{sec::eg-sensitivity-tau}中使用的数据。


\paragraph{Recommended reading}

 \citet{Rosenbaum::1983JRSSB} 和 \citet{imbens2003sensitivity} 是 sensitivity analysis 的两篇经典论文，不过它们涉及更复杂的步骤。
